\section{Introduction}
Large language model (LLM) agents use tools to retrieve information, invoke services, and modify external state on a user's behalf~\cite{react,toolemu}. These capabilities allow agents to carry out tasks in diverse software environments~\cite{agentbench}, but also give their decisions consequences for protected data and resources. Access control determines which resources an agent may access, which operations it may perform, and under what conditions it may act~\cite{progent,saga}. The central requirement is to provide sufficient authority for the task while keeping execution within the user's delegation and applicable organizational policies.

Authorizing an agent is difficult because both its course of action and the conditions governing permission can change during execution. Users typically specify a goal, while the agent selects concrete operations as it gathers information. Different paths to the same goal may access different resources, invoke different services, and produce different external effects; permission for one path need not cover another. New observations can lead the agent to revise its plan, delegation introduces additional principals, and delayed operations may outlive the approvals under which they were prepared. Authorization must therefore remain tied to the operations being performed and the conditions under which their effects occur. Successful task completion alone does not establish that the chosen execution stayed within the user's delegated authority.

Established security principles, such as least privilege~\cite{saltzer,privsep,capsicum} and complete mediation~\cite{saltzer,schneider,sasi}, have long guided the design of access control in software systems; applying these principles to agent workflows, however, requires accounting for decisions made by the agent itself. Least privilege limits authority to the permissions needed for execution, yet an agent can choose among plans with different resource and disclosure requirements. Complete mediation requires checking every protected operation, yet a single agent action may trigger effects across tools, delegated agents, and external services. Applying these principles therefore requires relating the agent's chosen plan to its authority requirements and tracing its actions to the boundaries where their effects can be controlled.

Recent research has advanced agent access control by addressing different aspects of authorization. Policy specification and derivation methods express constraints or infer them from task descriptions, organizational rules, and user preferences~\cite{conseca,agentspec,guardagent,shieldagent,miniscope,dataperms,drift}. Identity and delegation mechanisms establish who is acting, bind permissions to agents or tools, and constrain the scope of delegated access~\cite{delegation,saga,astra,aip,etdi}. Action and resource enforcement systems check operations at tool, application, browser, and operating-system boundaries~\cite{progent,ac4a,cellmate,agentbound,verisafe,csagent,pcas}. Information-flow control and supporting isolation mechanisms restrict how information reaches agent components, influences execution, or is released to other parties~\cite{camel,fides,ace,airgap,isolate,llmbda}. Together, these approaches provide complementary mechanisms for controlling agent authority throughout a workflow. Combining them requires preserving restrictions across interfaces: a policy limiting disclosure to an approved recipient, for example, must continue to constrain delegated agents and the tools through which they release information. Understanding these approaches requires examining what each mechanism enforces and how they compose.

In this paper, we aim to clarify the capabilities and limitations of existing approaches to agent access control and identify further directions. Our main contributions are as follows:
\begin{itemize}[leftmargin=*]
    \item \textbf{Surveying.} We survey and classify recent work on agent access control into the four mechanism families above. We compare their control objects, enforcement boundaries, and assumptions, and examine how these mechanisms jointly preserve authorization throughout an agent workflow.
    \item \textbf{Rethinking.} We draw on this analysis to rethink four aspects of agent access control: how plan selection shapes the authority an agent needs; when contextual information may inform authorization decisions; how restrictions survive translation across system interfaces; and whether authorization remains valid as workflows evolve.
    \item \textbf{Forecasting.} Building on these insights, we identify open challenges and directions for future research. We formulate concrete research questions and criteria for evaluating progress.
\end{itemize}
